Hem construït aquesta gràfica a partir de dades de freqüència recollides quan una font de so es movia acostant-se a nosaltres (velocitats positives) o allunyant-se'n (velocitats negatives), a velocitats diferents.

\figura[0.55]{fig1.png}

\begin{apartat}{1}
Com s'anomena el fenomen que hem estudiat en aquest experiment? La font de so s'acosta a nosaltres amb un moviment rectilini uniforme (MRU) a $100\ \mathrm{m\,s^{-1}}$ i ens sobrepassa. Quin canvi de freqüència (expressada en Hz) sentirem en el moment en què passi just pel nostre costat? La freqüència que sentirem augmentarà o disminuirà?
\end{apartat}

\begin{apartat}{1}
La taula següent mostra com disminueix la intensitat sonora quan ens situem a diferents distàncies d'un emissor puntual de so.

\begin{center}
\begin{tabular}{|c|c|c|c|c|c|c|c|}
\hline
\textit{Distància (m)} & 5,0 & 10,0 & 15,0 & 20,0 & 25,0 & 30,0 & 35,0\\
\hline
\textit{I (mW m$^{-2}$)} & 0,080 & 0,020 & 0,0089 & 0,0050 & 0,0032 & 0,0022 & 0,0016\\
\hline
\end{tabular}
\end{center}

Calculeu a quina distància, aproximadament, haurem d'estar perquè el nivell de sensació sonora sigui de 65~dB i calculeu la potència de la font sonora, suposant que emet igual en totes les direccions.
\end{apartat}

\blocetiquetat{Dada:}{Intensitat del llindar d'audició (0~dB), $I_0 = 1,00\times 10^{-12}\ \mathrm{W\,m^{-2}}$}
